\documentclass[11pt,a4paper]{article}

\usepackage[T1]{fontenc}
\usepackage[utf8]{inputenc}
\usepackage[english]{babel}
\usepackage{geometry}
\usepackage{listings}
\usepackage{xcolor}
\usepackage{booktabs}
\usepackage{longtable}
\usepackage{array}
\usepackage{hyperref}
\usepackage{fancyhdr}
\usepackage{titlesec}
\usepackage{enumitem}
\usepackage{graphicx}
\usepackage{amsmath}
\usepackage{parskip}

\geometry{
    top=25mm, bottom=25mm,
    left=25mm, right=25mm
}

% ---------------------------------------------------------------------------
% Colours
% ---------------------------------------------------------------------------
\definecolor{codebg}{RGB}{245,245,245}
\definecolor{codeframe}{RGB}{200,200,200}
\definecolor{kwcolor}{RGB}{0,0,180}
\definecolor{commentcolor}{RGB}{80,120,80}
\definecolor{stringcolor}{RGB}{160,32,32}
\definecolor{shellprompt}{RGB}{100,100,100}

% ---------------------------------------------------------------------------
% Listings style
% ---------------------------------------------------------------------------
\lstdefinestyle{shell}{
    backgroundcolor=\color{codebg},
    frame=single,
    framerule=0.4pt,
    rulecolor=\color{codeframe},
    basicstyle=\ttfamily\small,
    breaklines=true,
    breakatwhitespace=false,
    showstringspaces=false,
    tabsize=4,
    captionpos=b,
    xleftmargin=6pt,
    xrightmargin=6pt,
    aboveskip=8pt,
    belowskip=8pt,
}

\lstdefinestyle{vhdl}{
    backgroundcolor=\color{codebg},
    frame=single,
    framerule=0.4pt,
    rulecolor=\color{codeframe},
    basicstyle=\ttfamily\small,
    keywordstyle=\color{kwcolor}\bfseries,
    commentstyle=\color{commentcolor}\itshape,
    stringstyle=\color{stringcolor},
    morekeywords={library,use,entity,port,architecture,signal,process,
                  begin,end,if,then,else,elsif,rising_edge,std_logic,
                  std_logic_vector,unsigned,integer,constant,when,
                  inout,in,out,downto,others,not,and,or},
    breaklines=true,
    showstringspaces=false,
    tabsize=4,
    captionpos=b,
    xleftmargin=6pt,
    xrightmargin=6pt,
    aboveskip=8pt,
    belowskip=8pt,
}

\lstset{style=shell}

% ---------------------------------------------------------------------------
% Page style
% ---------------------------------------------------------------------------
\pagestyle{fancy}
\fancyhf{}
\fancyhead[L]{\small VERA RBL-XE -- FPGA Build \& Programming Guide}
\fancyhead[R]{\small Rev~2}
\fancyfoot[C]{\thepage}
\renewcommand{\headrulewidth}{0.4pt}

% ---------------------------------------------------------------------------
% Section formatting
% ---------------------------------------------------------------------------
\titleformat{\section}{\large\bfseries}{}{0em}{\thesection\quad}
\titleformat{\subsection}{\normalsize\bfseries}{}{0em}{\thesubsection\quad}

% ---------------------------------------------------------------------------
% Convenience macros
% ---------------------------------------------------------------------------
\newcommand{\cmd}[1]{\texttt{#1}}
\newcommand{\sig}[1]{\texttt{#1}}
\newcommand{\file}[1]{\texttt{#1}}
\newcommand{\hex}[1]{\texttt{\$#1}}

% ---------------------------------------------------------------------------
\begin{document}
% ---------------------------------------------------------------------------

\begin{titlepage}
    \centering
    \vspace*{30mm}
    {\Huge\bfseries VERA RBL-XE\\[4mm]
     FPGA Build \& Programming Guide\par}
    \vspace{10mm}
    {\large iCE40UP5K -- Rev~2\par}
    \vspace{8mm}
    {\normalsize
        6502 bus controller with RAMbo 256\,KB, PBI RAM/ROM,\\
        and external SRAM support\par}
    \vfill
    {\small Last updated: \today\par}
\end{titlepage}

\tableofcontents
\newpage

% ===========================================================================
\section{Overview}
% ===========================================================================

This document describes the complete workflow for building and programming
the FPGA firmware for the VERA RBL-XE module.  The firmware implements a
6502 bus controller on a Lattice iCE40UP5K, replacing the ESP32-S3 MCU
used in the previous design.

\subsection{Architecture summary}

\begin{itemize}[noitemsep]
    \item \textbf{PBI RAM} (\hex{D600}--\hex{D7FF}, 512\,B): stored in
          iCE40 BRAM (1 EBR block).
    \item \textbf{PBI ROM} (\hex{D800}--\hex{DFFF}, 2\,KB): stored in
          iCE40 BRAM (4 EBR blocks), pre-loaded at synthesis time from
          the assembled 6502 driver binary.
    \item \textbf{RAMbo} (\hex{4000}--\hex{7FFF}, 16\,KB window,
          256\,KB total): served by an external 512K$\times$8 SRAM
          (e.g.\ AS6C4008-55SIN).  The FPGA controls chip-enable,
          output-enable, write-enable, and the upper four address bits
          (bank), using only 7 new I/O pins.
    \item \textbf{PHI2 synchronisation}: PHI2 is treated as a data input
          and synchronised to the 25\,MHz FPGA clock via a two-stage
          flip-flop chain, producing a single-cycle \cmd{phi2\_rise}
          pulse used to trigger all registered logic.
\end{itemize}

\subsection{Memory map}

\begin{center}
\begin{tabular}{llll}
\toprule
\textbf{Range} & \textbf{Size} & \textbf{Function} & \textbf{Implementation} \\
\midrule
\hex{4000}--\hex{7FFF} & 16\,KB window & RAMbo (256\,KB banked) & External SRAM \\
\hex{D1FF}            & 1\,B          & VERA\_CS latch (write) & Register \\
\hex{D301}            & 1\,B          & PIA PORTB snoop        & Register \\
\hex{D303}            & 1\,B          & PIA PBCTL snoop        & Register \\
\hex{D600}--\hex{D7FF} & 512\,B        & PBI RAM                & BRAM (1 EBR) \\
\hex{D800}--\hex{DFFF} & 2\,KB         & PBI ROM                & BRAM (4 EBR) \\
\bottomrule
\end{tabular}
\end{center}

% ===========================================================================
\section{Prerequisites}
% ===========================================================================

\subsection{FPGA toolchain}

All tools below are open-source.  Install on Debian/Ubuntu:

\begin{lstlisting}[style=shell]
sudo apt install fpga-icestorm nextpnr-ice40 yosys ghdl
\end{lstlisting}

\textbf{Note:} the \cmd{ghdl-yosys-plugin} is required for VHDL synthesis
and is generally not available as a package.  Build it from source:

\begin{lstlisting}[style=shell]
git clone https://github.com/ghdl/ghdl-yosys-plugin
cd ghdl-yosys-plugin
make
sudo make install
\end{lstlisting}

\subsection{Flash programming tools}

\begin{lstlisting}[style=shell]
# openFPGALoader -- recommended, supports many programmers
sudo apt install openFPGALoader

# iceprog is already included in fpga-icestorm
\end{lstlisting}

\subsection{6502 assembler (cc65)}

The PBI ROM driver is written in 6502 assembly and built with the
cc65 toolchain.

\begin{lstlisting}[style=shell]
sudo apt install cc65
\end{lstlisting}

\subsection{Python 3}

Required to run \file{gen\_pbi\_rom\_pkg.py}, which converts the
assembled 6502 binary into a VHDL constant package.

\begin{lstlisting}[style=shell]
sudo apt install python3
\end{lstlisting}

\subsection{Tool roles at a glance}

\begin{center}
\begin{tabular}{llll}
\toprule
\textbf{Tool} & \textbf{Input} & \textbf{Output} & \textbf{Purpose} \\
\midrule
\cmd{ca65} / \cmd{ld65}    & \file{.s} sources    & \file{.bin}      & 6502 assembly \\
\cmd{gen\_pbi\_rom\_pkg.py} & \file{.bin}          & \file{.vhd}      & ROM initialisation \\
\cmd{ghdl} (plugin)         & VHDL sources         & RTLIL            & VHDL elaboration \\
\cmd{yosys}                 & RTLIL                & \file{.json}     & Synthesis \\
\cmd{nextpnr-ice40}         & \file{.json + .pcf}  & \file{.asc}      & Place \& route \\
\cmd{icepack}               & \file{.asc}          & \file{.bin}      & Bitstream packing \\
\cmd{openFPGALoader}        & \file{.bin}          & SPI flash        & Programming \\
\bottomrule
\end{tabular}
\end{center}

% ===========================================================================
\section{Build Flow}
% ===========================================================================

\subsection{Step 1 -- Assemble the 6502 ROM driver}

The PBI ROM at \hex{D800}--\hex{DFFF} contains a 2\,KB 6502 driver
assembled from \file{firmware/MCU/6502/src/vera\_pbi\_handler.s}.

\begin{lstlisting}[style=shell,caption={Assemble the 6502 driver}]
cd firmware/MCU/6502
make
\end{lstlisting}

This single command runs three sub-steps:

\begin{enumerate}[noitemsep]
    \item \textbf{Assemble}: \cmd{ca65} compiles
          \file{vera\_pbi\_handler.s} into \file{vera\_pbi\_handler.o}.
    \item \textbf{Link}: \cmd{ld65} links with \file{pbi-driver.ld}
          to produce \file{vera\_pbi\_handler.bin}.
    \item \textbf{Generate VHDL ROM package}: \file{gen\_pbi\_rom\_pkg.py}
          reads the binary and writes
          \file{../../FPGA/pbi\_rom\_pkg.vhd} with the
          \cmd{PBI\_ROM\_INIT} constant.
\end{enumerate}

To regenerate only the VHDL package when the binary already exists:

\begin{lstlisting}[style=shell]
make fpga-rom
\end{lstlisting}

\subsection{Step 2 -- Synthesis}

Synthesis converts the VHDL sources into a technology-mapped netlist
for the iCE40 family.

\begin{lstlisting}[style=shell,caption={Synthesis with Yosys + GHDL plugin}]
cd firmware/FPGA
yosys -m ghdl -p \
    "ghdl --std=08 custom_types.vhd pbi_rom_pkg.vhd C6502_Interface.vhd \
          -e C6502_Interface; \
     synth_ice40 -top C6502_Interface -json C6502_Interface.json"
\end{lstlisting}

The \cmd{-m ghdl} flag loads the GHDL plugin into Yosys, which handles
VHDL analysis and elaboration.  The \cmd{synth\_ice40} pass maps the
resulting netlist to iCE40 primitives (LUTs, FFs, BRAM, SB\_IO, etc.)
and writes a JSON netlist.

On success, Yosys prints a resource utilisation summary.  Expected ranges:

\begin{center}
\begin{tabular}{lr}
\toprule
\textbf{Resource} & \textbf{Expected usage} \\
\midrule
SB\_LUT4 (logic cells) & 80--180 \\
SB\_FF (flip-flops)    & 30--60  \\
SB\_RAM40\_4K (EBR)    & 5       \\
SB\_IO (I/O buffers)   & 40      \\
\bottomrule
\end{tabular}
\end{center}

\subsection{Step 3 -- Place \& Route}

Place \& route assigns logic cells and routes signals to physical
iCE40 resources, guided by the pin constraints file.

\begin{lstlisting}[style=shell,caption={Place \& route with nextpnr-ice40}]
nextpnr-ice40 --up5k --package sg48 \
    --json C6502_Interface.json \
    --pcf board-pin-mapping.pcf \
    --asc C6502_Interface.asc
\end{lstlisting}

\textbf{Important:} all pin numbers in \file{board-pin-mapping.pcf} are
placeholders.  Edit the file to match your actual schematic before this
step (see Section~\ref{sec:pins}).

nextpnr reports timing after routing.  The design must meet timing at
25\,MHz (the system clock).  The critical path is the PHI2 synchroniser
chain, which is combinational; the path from \cmd{phi2\_rise} to any
registered output is unconstrained by PHI2 frequency.

\subsection{Step 4 -- Pack the bitstream}

\cmd{icepack} converts the ASCII intermediate format produced by
nextpnr into the binary bitstream that is written to the SPI flash.

\begin{lstlisting}[style=shell,caption={Pack binary bitstream}]
icepack C6502_Interface.asc C6502_Interface.bin
\end{lstlisting}

\cmd{C6502\_Interface.bin} is the file that gets programmed into the
SPI flash.  Its size is typically around 100--200\,KB depending on
the device configuration.

\subsection{Step 5 -- Program the SPI flash}

The iCE40UP5K loads its configuration from an external SPI NOR flash
at power-up.  The programmer writes the bitstream into that flash.

\subsubsection{Using openFPGALoader (recommended)}

\begin{lstlisting}[style=shell,caption={Program with openFPGALoader}]
openFPGALoader -b ice40_generic C6502_Interface.bin
\end{lstlisting}

Select the programmer cable explicitly if needed:

\begin{lstlisting}[style=shell]
openFPGALoader --cable ft2232    C6502_Interface.bin   # FTDI FT2232H
openFPGALoader --cable ft232     C6502_Interface.bin   # FTDI FT232H
openFPGALoader --cable tigard    C6502_Interface.bin   # Tigard
openFPGALoader --cable raspberrypi C6502_Interface.bin # Raspberry Pi GPIO
\end{lstlisting}

\subsubsection{Using iceprog (FTDI FT2232H)}

\begin{lstlisting}[style=shell,caption={Program with iceprog}]
iceprog C6502_Interface.bin

# Erase + write (useful if previous write was partial)
iceprog -e 1048576 C6502_Interface.bin
\end{lstlisting}

\subsection{One-command build and flash}

The \file{firmware/FPGA/Makefile} wraps all steps:

\begin{lstlisting}[style=shell]
cd firmware/FPGA
make             # synthesise + P&R + pack
make flash       # write with openFPGALoader
make flash-iceprog   # write with iceprog
make clean       # remove generated files
\end{lstlisting}

% ===========================================================================
\section{Programmer Wiring}
% ===========================================================================

The SPI flash is shared between the FPGA (configuration bus at boot)
and the programmer (writes the bitstream).  During programming, the
FPGA is held in reset by asserting \sig{CRESET\_B} low.

\begin{center}
\begin{tabular}{lll}
\toprule
\textbf{Programmer signal} & \textbf{iCE40UP5K pin} & \textbf{SPI flash pin} \\
\midrule
MOSI  & SPI\_SI   & DI  \\
MISO  & SPI\_SO   & DO  \\
SCK   & SPI\_SCK  & CLK \\
CS\#  & --        & CS\# \\
CRST\# & CRESET\_B & -- \\
CDONE & CDONE     & -- \\
\bottomrule
\end{tabular}
\end{center}

Recommended SPI flash: Winbond W25Q16 (2\,MB) or W25Q32 (4\,MB).

% ===========================================================================
\section{Pin Assignment}
\label{sec:pins}
% ===========================================================================

Edit \file{board-pin-mapping.pcf} and replace every pin number with the
actual pad number from your schematic before running nextpnr-ice40.

The \sig{CLK} signal must be assigned to a \textbf{Global Buffer (GB)
input} pin.  Valid GB pins on the iCE40UP5K SG48 package are:
\textbf{20, 28, 31, 38, 41, 43, 48}.

\subsection{Complete pin table}

\begin{center}
\small
\begin{longtable}{llllp{5.5cm}}
\toprule
\textbf{Signal} & \textbf{Dir} & \textbf{PCF pin} & \textbf{Std} &
\textbf{Description} \\
\midrule
\endfirsthead
\toprule
\textbf{Signal} & \textbf{Dir} & \textbf{PCF pin} & \textbf{Std} &
\textbf{Description} \\
\midrule
\endhead
\midrule
\multicolumn{5}{r}{\small continued on next page}\\
\endfoot
\bottomrule
\endlastfoot

% --- Clock ---
\sig{CLK}           & In    & 35* & LVCMOS33 &
    25\,MHz system clock. Must land on a GB pin. \\[2pt]

% --- 6502 address bus ---
\sig{A[0]}          & In    & 40* & LVCMOS33 & 6502 address bit 0 \\
\sig{A[1]}          & In    & 41* & LVCMOS33 & 6502 address bit 1 \\
\sig{A[2]}          & In    & 42* & LVCMOS33 & 6502 address bit 2 \\
\sig{A[3]}          & In    & 43* & LVCMOS33 & 6502 address bit 3 \\
\sig{A[4]}          & In    & 44* & LVCMOS33 & 6502 address bit 4 \\
\sig{A[5]}          & In    & 45* & LVCMOS33 & 6502 address bit 5 \\
\sig{A[6]}          & In    & 46* & LVCMOS33 & 6502 address bit 6 \\
\sig{A[7]}          & In    & 47* & LVCMOS33 & 6502 address bit 7 \\
\sig{A[8]}          & In    & 48* & LVCMOS33 & 6502 address bit 8 \\
\sig{A[9]}          & In    &  1* & LVCMOS33 & 6502 address bit 9 \\
\sig{A[10]}         & In    &  2* & LVCMOS33 & 6502 address bit 10 \\
\sig{A[11]}         & In    &  3* & LVCMOS33 & 6502 address bit 11 \\
\sig{A[12]}         & In    &  4* & LVCMOS33 & 6502 address bit 12 \\
\sig{A[13]}         & In    &  7* & LVCMOS33 &
    6502 address bit 13. Also wired to SRAM A[13] on PCB. \\
\sig{A[14]}         & In    &  9* & LVCMOS33 & 6502 address bit 14 \\
\sig{A[15]}         & In    & 10* & LVCMOS33 & 6502 address bit 15 \\[2pt]

% --- 6502 data bus ---
\sig{D[0]}          & InOut & 11* & LVCMOS33 &
    6502 data bit 0. Shared net with SRAM DQ[0]. \\
\sig{D[1]}          & InOut & 12* & LVCMOS33 &
    6502 data bit 1. Shared net with SRAM DQ[1]. \\
\sig{D[2]}          & InOut & 13* & LVCMOS33 &
    6502 data bit 2. Shared net with SRAM DQ[2]. \\
\sig{D[3]}          & InOut & 14* & LVCMOS33 &
    6502 data bit 3. Shared net with SRAM DQ[3]. \\
\sig{D[4]}          & InOut & 15* & LVCMOS33 &
    6502 data bit 4. Shared net with SRAM DQ[4]. \\
\sig{D[5]}          & InOut & 16* & LVCMOS33 &
    6502 data bit 5. Shared net with SRAM DQ[5]. \\
\sig{D[6]}          & InOut & 17* & LVCMOS33 &
    6502 data bit 6. Shared net with SRAM DQ[6]. \\
\sig{D[7]}          & InOut & 18* & LVCMOS33 &
    6502 data bit 7. Shared net with SRAM DQ[7]. \\[2pt]

% --- 6502 control ---
\sig{PHI2}          & In    & 34* & LVCMOS33 &
    6502 phase-2 clock. Synchronised internally; not used as FPGA clock. \\
\sig{RNW\_}          & In    & 36* & LVCMOS33 &
    Read/Not-Write. High = read, Low = write. \\[2pt]

% --- PBI control ---
\sig{VERA\_CS}      & Out   & 32* & LVCMOS33 &
    VERA chip-select latch output (active low). \\
\sig{MPD}           & Out   & 31* & LVCMOS33 &
    Machine Program Data. Asserted low for \hex{D800}--\hex{DFFF}
    reads; blocks Atari from driving D[7:0]. \\
\sig{EXTSEL}        & Out   & 30* & LVCMOS33 &
    External Select. Asserted low for PBI RAM (\hex{D600}--\hex{D7FF})
    and RAMbo (\hex{4000}--\hex{7FFF}) accesses; blocks Atari internal
    RAM from responding. \\[2pt]

% --- DIP switches ---
\sig{DIP\_SEL[0]}   & In    & 27* & LVCMOS33 &
    DIP switch bit 0. Selects which D[n] bit controls VERA\_CS latch. \\
\sig{DIP\_SEL[1]}   & In    & 28* & LVCMOS33 & DIP switch bit 1. \\
\sig{DIP\_SEL[2]}   & In    & 29* & LVCMOS33 & DIP switch bit 2. \\[2pt]

% --- External SRAM ---
\sig{SRAM\_A\_BANK[0]} & Out & 19* & LVCMOS33 &
    SRAM address bit 14 (bank LSB). \\
\sig{SRAM\_A\_BANK[1]} & Out & 20* & LVCMOS33 &
    SRAM address bit 15. \\
\sig{SRAM\_A\_BANK[2]} & Out & 21* & LVCMOS33 &
    SRAM address bit 16. \\
\sig{SRAM\_A\_BANK[3]} & Out & 22* & LVCMOS33 &
    SRAM address bit 17 (bank MSB). \\
\sig{SRAM\_CE\_N}   & Out   & 23* & LVCMOS33 &
    SRAM chip enable (active low). Asserted during PHI2-high when
    address is in \hex{4000}--\hex{7FFF} and RAMbo is enabled. \\
\sig{SRAM\_OE\_N}   & Out   & 24* & LVCMOS33 &
    SRAM output enable (active low). Asserted on read cycles only. \\
\sig{SRAM\_WE\_N}   & Out   & 25* & LVCMOS33 &
    SRAM write enable (active low). Asserted on write cycles only.
    Pulse width equals PHI2-high time ($\approx$279\,ns at 1.79\,MHz). \\

\end{longtable}
\end{center}

\textit{* Pin numbers are placeholders.  Replace with actual pad numbers
from your board schematic.}

\subsection{External SRAM address wiring}

The SRAM address bus is split between PCB traces (lower 14 bits) and
FPGA outputs (upper 4 bank bits):

\begin{center}
\begin{tabular}{lll}
\toprule
\textbf{SRAM address bit} & \textbf{Source} & \textbf{Note} \\
\midrule
A[13:0]  & 6502 A[13:0] (PCB net) & Direct trace, no FPGA I/O pin \\
A[14]    & FPGA \sig{SRAM\_A\_BANK[0]} & Bank bit 0 = PORTB[2] \\
A[15]    & FPGA \sig{SRAM\_A\_BANK[1]} & Bank bit 1 = PORTB[3] \\
A[16]    & FPGA \sig{SRAM\_A\_BANK[2]} & Bank bit 2 = PORTB[5] \\
A[17]    & FPGA \sig{SRAM\_A\_BANK[3]} & Bank bit 3 = PORTB[6] \\
A[18]    & GND (PCB)               & Selects lower 256\,KB of 512\,KB chip \\
\bottomrule
\end{tabular}
\end{center}

\subsection{RAMbo bank formula}

The FPGA snoops writes to the Atari PIA to derive the active bank:

\begin{lstlisting}[style=shell]
PORTB at $D301  -- bank bits
PBCTL at $D303  -- bit 2 enables RAMbo

bank[1:0] = PORTB[3:2]
bank[3:2] = PORTB[6:5]
=> SRAM A[17:14] = { PORTB[6], PORTB[5], PORTB[3], PORTB[2] }
\end{lstlisting}

RAMbo is active only when \sig{PBCTL[2] = 1}.

% ===========================================================================
\section{Timing}
% ===========================================================================

\subsection{Read cycle (SRAM)}

\begin{center}
\begin{tabular}{lrl}
\toprule
\textbf{Event} & \textbf{Time from PHI2 rise} & \textbf{Note} \\
\midrule
Address valid (6502 tADS) & $-$40\,ns & Before PHI2 rise \\
SRAM\_CE\# asserted       & $\approx$0\,ns   & Combinational gate on PHI2 \\
SRAM data valid (tAA=55\,ns) & $+$55\,ns & AS6C4008-55 \\
6502 data setup deadline  & $+$249\,ns & 279\,ns $-$ 30\,ns tDSR \\
\textbf{Read margin}      & \textbf{$+$194\,ns} & Adequate \\
\bottomrule
\end{tabular}
\end{center}

\subsection{Write cycle (SRAM)}

WE\# is asserted combinationally when PHI2 is high and the address is
in the RAMbo window.  The pulse width equals the PHI2-high time:

\[
t_{WC} = t_{\text{PHI2-high}} \approx 279\,\text{ns}
\quad \gg \quad
t_{WC(\text{SRAM})} = 55\,\text{ns}
\]

\subsection{PBI RAM / ROM (BRAM)}

BRAM reads are triggered on \cmd{phi2\_rise} (one clock cycle after the
real PHI2 edge, $\approx$40\,ns delay).  The BRAM output is registered
and valid within one additional clock cycle ($\approx$40\,ns), leaving
$\approx$199\,ns of margin before the 6502 samples the bus.

% ===========================================================================
\section{Source File Reference}
% ===========================================================================

\begin{center}
\begin{tabular}{lp{9cm}}
\toprule
\textbf{File} & \textbf{Description} \\
\midrule
\file{C6502\_Interface.vhd} &
    Top-level VHDL entity (Rev~2). Contains all bus logic, BRAM
    instances, and SRAM control. \\
\file{custom\_types.vhd} &
    Package defining \cmd{rom\_array} (2048$\times$8) and
    \cmd{ram\_512\_array} (512$\times$8) array types. \\
\file{pbi\_rom\_pkg.vhd} &
    Auto-generated package containing the \cmd{PBI\_ROM\_INIT}
    constant (2048 bytes).  Do not edit manually. \\
\file{gen\_pbi\_rom\_pkg.py} &
    Python~3 script that converts a 6502 \file{.bin} file into
    \file{pbi\_rom\_pkg.vhd}.  Invoked by the MCU Makefile. \\
\file{board-pin-mapping.pcf} &
    Pin constraints for iCE40UP5K SG48.
    All pin numbers are placeholders -- edit before use. \\
\file{Makefile} &
    Targets: \cmd{all}, \cmd{flash}, \cmd{flash-iceprog}, \cmd{clean}. \\
\bottomrule
\end{tabular}
\end{center}

\end{document}
